\section{Introduction and Motivation}
\label{sec:introduction}

Knowledge graphs organize entities and typed relations into networks that
support multi-hop queries. In biomedicine, resources such as PrimeKG link
drugs, diseases, and proteins so that associations can be traced across
connected biological facts \cite{chandak2023primekg}. When links or entities
are lost, however, a graph can remain almost fully connected while failing
the queries it is meant to answer. On our 5,000-node PrimeKG subgraph,
randomly removing 50\% of the edges leaves 98.2\% of the original nodes in
the largest connected component, yet only 43.3\% of 200 benchmark queries
retain a valid two-hop path---a 54.9-percentage-point gap measured on the
same damage states and averaged over ten seeds (Section~\ref{sec:functional}).

This divergence arises because global connectivity and typed retrieval
impose different structural requirements. Classical network-robustness
studies evaluate damage by tracking the largest connected component under
random failures and targeted removals \cite{albert2000error,callaway2000network}.
That statistic treats all edges as interchangeable and permits paths of
arbitrary length. A structured query, by contrast, requires specific
relation types within a bounded hop limit. In our benchmark, a two-hop
drug--disease query succeeds only if a shared protein retains both its
drug--protein and disease--protein links. Even after those links disappear,
dense protein--protein interactions can still keep both endpoints inside
the giant component through longer routes. Global connectivity can
therefore remain high after short, typed association paths have largely
collapsed.

Distinguishing structural connectivity from retrieval usability matters for
both evaluation and repair. Because the largest connected component is a
coarse network-level statistic, a system that relies on typed paths can
lose more than half of its usable query support while connectivity barely
changes. Measuring retrieval directly is therefore necessary to assess
functional degradation. The same distinction reshapes network restoration:
when only a limited budget of lost edges can be repaired, the links that
best rebuild global connectivity need not be the links that best recover
failed queries, and edge priorities may shift when a task can use longer
paths.

Motivated by this gap, we investigate how random failures and targeted
removals affect global connectivity compared with typed-path retrieval on
the same damaged graphs. We then examine how path redundancy---multiple
shared intermediate proteins---and the allowed number of hops buffer
retrieval against edge loss. Finally, under a fixed repair budget, we ask
which deleted edges best recover global connectivity versus held-out query
success, and how those priorities change with the query horizon and budget.

Using a biomedical benchmark derived from PrimeKG \cite{chandak2023primekg},
we combine measurement with theory. We pair structural robustness analysis
with direct evaluation of typed drug--disease retrieval on matched damage
states (Sections~\ref{sec:structural} and~\ref{sec:functional}). We derive
zero-parameter predictions for both kinds of robustness and test where they
hold and where they fail (Sections~\ref{sec:theory} and~\ref{sec:prediction}).
Finally, we compare limited-budget restoration strategies on connectivity and
on held-out query recovery (Section~\ref{sec:restoration}). Our main findings
are the following.
\begin{enumerate}[leftmargin=*,itemsep=2pt]
\item \textbf{Connectivity and retrieval are different kinds of quantity.}
The giant component is a collective property whose percolation threshold, for
this graph, lies beyond 99\% edge loss. Two-hop retrieval is a local
reliability problem. Its success rate is a polynomial,
$R_2=1-G_m(1-p^2)$, determined by the number $m$ of independent routes per
query. This formula has no fitted parameter and matches the long-run
simulations within 0.001.
\item \textbf{The query horizon decides which redundancy counts.} Allowing
longer paths moves the bottleneck from shared proteins to the drugs' own target
links. For every query, the drug--disease minimum cut equals the smaller
endpoint degree, and the expected long-range percolation collapse does not
appear.
\item \textbf{Targeted attacks break the independence that makes retrieval
predictable.} Query routes pass through the same hubs. Removing the top 5\% of
nodes by degree lowers the giant component only 6.6 percentage points more
than random removal, but lowers retrieval 56.5 points more. Degree-preserving
null models trace the excess fragility to PrimeKG's habit of hanging
low-degree entities on hubs.
\item \textbf{Repair priorities follow the task.} At two hops, restoring only
drug and disease links is best, as the formula predicts within one seed SD.
At three hops, structural repair becomes competitive.
\end{enumerate}
